% V2 PDF物理页73--74 原书思考题14.1--14.11；图14.13等待专家校核。
\section*{第14章\quad 静电场中的导体：思考题}
\paragraph{思考题 14.1}
具有各种几何形状的导体置于静电平衡后，是否只有球形导体内部电场才为零？

\paragraph{思考题 14.2}
总电荷为Q的孤立导体球壳中心放入点电荷q。若原球壳电势为φ₀，能否把插入后的电势处处写成点电荷势加原球壳电势φ₀？

\paragraph{思考题 14.3}
把带正电的孤立导体金属球加入一批电子，球的质量如何改变？电能变化与电子质量怎样区分？

\paragraph{思考题 14.4}
球形导体壳的中心有点电荷时内外表面电荷分布是否均匀？若点电荷偏离球心会怎样？

\paragraph{思考题 14.5}
带电物体放在导体空腔附近，为什么导体内的静电屏蔽可让外部电场不影响封闭空腔中的静电场？

\paragraph{思考题 14.6}
导体表面某点附近有面电荷密度(sigma)，为何其表面外侧总电场 (E_n=sigma/epsilon_0)？若移近别的带电物体，电场与局部(sigma)关系是否仍成立？

\paragraph{思考题 14.7}
空间中两导体各有静电荷，证明至少有一个导体表面某点所带面电荷是正的或至少不是负的，讨论两者总电荷的条件。

\paragraph{思考题 14.8}
为什么孤立无穷均匀带电面两侧各得场强(sigma/(2epsilon_0))，但带电导体的表面外侧为(sigma/epsilon_0)？

\paragraph{思考题 14.9}
两个平行导体大平板带电后，相对的表面是否一定有等量异号电荷？讨论边缘场和外表面场。

\paragraph{思考题 14.10}
将原本放在导体球壳中心的点电荷q移到球壳外部，导体球本身的电势怎样变化？区分接地与孤立导体。

\paragraph{思考题 14.11}
两组导线相连的金属球A/B初始无净电荷，另有外部带异号电的C/D球靠近并产生感应。若将C/D用导线相连，A/B上已感应电荷的符号与数量是否可能改变？
